\chapter{Análisis y Diseño}
	\label{cap:analisis_diseno}
	Este capítulo documenta el análisis y el diseño de la plataforma tal como quedó definida al
	cierre de esta etapa del Trabajo Terminal. Su propósito es dejar constancia de \emph{qué} se
	construyó y, sobre todo, de \emph{por qué} se construyó de esa manera: los requerimientos que
	se levantaron con la Defensoría, los actores que intervienen, los casos de uso que delimitan
	el alcance funcional, la arquitectura que soporta la solución y el modelo de datos que la
	sustenta. La construcción concreta de cada componente ---el código, el despliegue y las
	pruebas--- se documenta en el Capítulo~\ref{cap:implementacion}, y el recorrido funcional de
	cada módulo entregado en el Capítulo~\ref{cap:modulos}.
	
	Los artefactos que aquí se presentan siguen el orden en que fueron producidos: primero el
	levantamiento de requerimientos con el cliente, del que se derivaron los casos de uso; después
	el modelo de estados del expediente, que ordena el flujo de trabajo; sobre esa base la
	arquitectura de componentes y el modelo de datos; después el diseño de los componentes de
	inteligencia artificial y de la interfaz de usuario; y, a lo largo de todo el capítulo, las
	decisiones de diseño con las alternativas que se descartaron en cada punto. Este orden no es una convención editorial:
	cada artefacto es insumo del siguiente, y esa dependencia es la que sostiene la coherencia
	técnica del sistema.
	
